\lecture{Neumann assembly and the first error estimate}{Lecture 6 --- Friday, September 4, 2026}{Math 6660 --- \texttt{6th.pdf}}

\begin{roadmap}[Lecture 6]
\roadq{How wrong is the FEM answer, and measured how?}
\begin{roadlist}
\item Assemble the Neumann problem explicitly first --- the boundary rows are
  where the natural BC shows up.
\item Then error analysis begins. Two \textbf{norms} matter: the energy norm
  (strain energy stored in the error) and the $L^{2}$ norm (plain size of the
  error).
\item First \textbf{a priori} estimate: the FEM error is bounded by the
  \emph{interpolation} error --- by how well the mesh could represent $u$ at
  best, regardless of the method.
\end{roadlist}
\payoff{Convergence studies: knowing what slope to expect before you run one.}
\end{roadmap}

% ==================================================================
\subsection*{Setup, restated}

\begin{equation}
  -(au')'=f \ \text{ in }(0,1),
  \qquad
  \underbrace{u(0)=0}_{\text{Dirichlet}},
  \qquad
  \underbrace{a(1)u'=g_{1}}_{\text{Neumann}}
\end{equation}
\begin{equation}
  \Vsp=\Bigl\{v:\textstyle\int\bigl[v^{2}+(v')^{2}\bigr]\dx<\infty,\ \
  \underbrace{v(0)=0}_{\text{strongly enforced}}\Bigr\}
\end{equation}

\paragraph{CG(1) method.}
\begin{equation}
  U(x)=\sum_{j=1}^{M+1}\xi_{j}\phih_{j}(x)
  \qquad
  \text{\small trial solution; the $+1$ is because the last node is unknown}
\end{equation}
\begin{equation}\label{eq:neuweak}
  \int_{0}^{1} a U'v'\dx = \int_{0}^{1} f v\dx
  + \underbrace{g_{1}v(1)}_{\text{Neumann --- weakly enforced}},
  \qquad v=\phih_{i},\ \ i=1,2,\dots,M+1 .
\end{equation}

% ==================================================================
\subsection*{Example: $a=1$, uniform mesh $h_{i}\equiv h$}

\begin{equation}
  \sum_{j=1}^{M+1}\xi_{j}\int_{0}^{1}\phih_{i}'\phih_{j}'\dx
  = \int_{0}^{1} f\phih_{i}\dx + g_{1}\phih_{i}(1)
\end{equation}

The new entries, all on the last row and column:
\begin{align}
  A_{(M,M+1)} &= \int_{x_{M}}^{x_{M+1}}\phih_{M+1}'\,\phih_{M}'\dx
    = -\frac{1}{h^{2}}\int_{x_{M}}^{x_{M+1}}\dx = -\frac{1}{h},\\
  A_{(M+1,M)} &= -\frac{1}{h}
    \qquad\Longrightarrow\quad \text{symmetric matrix},\\
  A_{(M+1,M+1)} &= \int_{x_{M}}^{x_{M+1}}\bigl(\phih_{M+1}'\bigr)^{2}\dx
    = \frac{1}{h}.
\end{align}

So $\underline{\underline{A}}\,\underline{\xi}=\underline{b}$ with
$\underline{\xi}=(\xi_{1},\xi_{2},\xi_{3},\dots,\xi_{M+1})$ and
\begin{equation}\label{eq:Aneu}
  A=\frac{1}{h}
  \begin{bmatrix}
    2 & -1 &        &    &    \\
   -1 &  2 & -1     &    &    \\
      & \ddots & \ddots & \ddots & \\
      &    & -1 & 2 & -1\\
      &    &    & -1 & \boxed{1}
  \end{bmatrix},
  \qquad
  A\in\R^{(M+1)\times(M+1)},\ \text{symmetric, positive definite.}
\end{equation}
\begin{equation}
  b_{i}=\int_{x_{i-1}}^{x_{i+1}} f\phih_{i}\dx \ \ (i=1,\dots,M),
  \qquad
  b_{M+1}=\int_{x_{M}}^{x_{M+1}} f\phih_{M+1}\dx + g_{1}.
\end{equation}

\begin{intuition}[the last diagonal entry is $1$, not $2$]
Every interior node has two elements attached, contributing $1/h$ each: hence
$2/h$. The end node $x_{M+1}$ has only \emph{one} element, so it gets $1/h$.

The $2,-1$ pattern was never a rule --- it is just ``sum of the elements meeting
at this node.'' This is exactly element-by-element assembly: each element
contributes $\frac{1}{h}\begin{psmallmatrix}1&-1\\-1&1\end{psmallmatrix}$ to its
two nodes, and the boundary node collects one copy instead of two.

Note also $g_{1}$ enters $b_{M+1}$ alone, as a plain \emph{nodal force} --- the
end traction assembled to the end node. No mystery.
\end{intuition}

\medskip
\noindent\textbf{Note: for the 2-pt BVP.}
\begin{enumerate}[label=\textcircled{\scriptsize\arabic*},leftmargin=2.2em,itemsep=2pt]
  \item Dirichlet BCs are \textbf{strongly} enforced \ (essential).
  \item Neumann or mixed BCs are \textbf{weakly} enforced \ (natural BCs).
\end{enumerate}

\begin{dictionary}[the two BC types]
essential / Dirichlet / strong & prescribed displacement; built into $\Vsp$; removes a DOF; exact\\
natural / Neumann / weak & applied traction; appears in $\Fvec$; adds no constraint; approximate\\
\end{dictionary}

% ==================================================================
\subsection*{Error analysis}

Two kinds of error estimates:
\begin{enumerate}[label=\textcircled{\scriptsize\arabic*},leftmargin=2.2em,itemsep=4pt]
  \item \textbf{a priori}: error estimate in terms of derivatives of the exact
        solution, \eg
        \[
          \underbrace{\norm{u}}_{}\;\dots\;
          \norm{\underbrace{u}_{\text{exact}}-\underbrace{U}_{\text{FEA}}}
          \;\le\; \underbrace{C}_{\text{constant}}\,
          h^{\,\underbrace{p}_{\mathclap{\text{order of shape function}}}}\,
          \norm{u^{(\ell)}}
          \qquad \ell^{\text{th}}\text{ derivative of } u .
        \]
  \item \textbf{a posteriori}: ``error estimates in terms of a residual'', \eg
        $\mathcal{L}(U)-f$. Useful for calculating error \emph{during} the
        simulation $\Rightarrow$ adaptive error control (adaptive mesh
        refinement).
\end{enumerate}

\begin{intuition}[before vs.\ after]
A priori uses $u$, which you do not have --- it predicts the \emph{rate}
($h^{p}$), so it tells you what refinement buys before you run anything.
A posteriori uses $U$, which you do have --- it says \emph{where} the error is,
which is what drives adaptive remeshing.
\end{intuition}

% ------------------------------------------------------------------
\subsection*{Norms}

Let $e(x)\defeq u(x)-U(x)$ be the error.
\begin{align}
  \text{``Energy'' norm:}  &\qquad \norm{v}_{E} = \sqrt{\int_{0}^{1} a(x)\bigl[v'(x)\bigr]^{2}\dx}\\
  \text{weighted $L_{2}$-norm:} &\qquad \norm{w}_{a} = \sqrt{\int_{0}^{1} a(x)\bigl[w(x)\bigr]^{2}\dx}
\end{align}
\begin{equation}\label{eq:normlink}
  \textbf{Note: } \norm{v}_{E}=\norm{v'}_{a}.
\end{equation}

\begin{mechanics}
$\norm{v}_{E}^{2}=\int a(v')^{2}=2\times$ strain energy. Measuring error in the
energy norm means measuring the strain energy of the error --- \ie how much
stored energy your mesh is missing. It weights \emph{strain} error, not
displacement error, which is why stress converges more slowly than displacement.
\end{mechanics}

\paragraph{Cauchy inequalities.}
\begin{align}
  \left|\int_{0}^{1} a v'w'\dx\right| &\le \norm{v'}_{a}\norm{w'}_{a}
     = \norm{v}_{E}\norm{w}_{E}, \label{eq:cauchyE}\\
  \left|\int_{0}^{1} v w\dx\right| &\le \norm{v}_{a}\norm{w}_{a^{-1}}.
\end{align}

% ------------------------------------------------------------------
\subsection*{An a priori error estimate}

CG(1) for homogeneous Dirichlet BCs: $-(au')'=f$ in $(0,1)$, $u(0)=u(1)=0$.

\paragraph{Error calculation.}
\begin{align}
  \int_{0}^{1} a U'v'\dx &= \int_{0}^{1} f v\dx \qquad \forall\, v\in\Vh\subset\Vsp, \\
  \int_{0}^{1} a u'v'\dx &= \int_{0}^{1} f v\dx \qquad \forall\, v\in\Vsp .
\end{align}
Now subtract:
\begin{equation}\label{eq:GO}
  \boxed{\;\int_{0}^{1} a\,(u-U)'\,v'\dx = 0 \qquad \forall\, v\in\Vh\;}
  \qquad \text{\textbf{Galerkin orthogonality}}
\end{equation}

\paragraph{Energy norm error estimate.} Using \eqref{eq:normlink},
\begin{equation}
  \norm{u-U}_{E}^{2} = \norm{(u-U)'}_{a}^{2}
  = \int_{0}^{1} a\,(u-U)'(u-U)'\dx .
\end{equation}
For \emph{any} $v\in\Vh$, insert $-v+v$:
\begin{align}
  \norm{u-U}_{E}^{2}
  &= \int_{0}^{1} a\,(u-U)'\,(u-v+v-U)'\dx \notag\\
  &= \int_{0}^{1} a\,(u-U)'(u-v)'\dx
   + \underbrace{\int_{0}^{1} a\,(u-U)'(v-U)'\dx}_{=\;0 \text{ by Galerkin orthogonality}} .
\end{align}

\medskip
\noindent\textbf{Note.} $u-v\in\Vsp$, but $v-U\in\Vh$ --- which is what lets
\eqref{eq:GO} kill the second term.

\begin{pitfall}[the trick is choosing which difference lies in $\Vh$]
$u-U$ is \emph{not} in $\Vh$ ($u$ is not a piecewise linear). But $v-U$
\emph{is}, since both $v$ and $U$ are. Galerkin orthogonality only annihilates
against members of $\Vh$, so the split $u-U=(u-v)+(v-U)$ is engineered to put a
$\Vh$ member in the second slot.
\end{pitfall}

\medskip
So, for any $v\in\Vh$,
\begin{equation}
  \norm{u-U}_{E}^{2} = \int_{0}^{1} a\,(u-U)'(u-v)'\dx .
\end{equation}
Using Cauchy \eqref{eq:cauchyE},
\begin{equation}
  \int_{0}^{1} a\,(u-U)'(u-v)'\dx \le \norm{(u-U)'}_{a}\,\norm{(u-v)'}_{a}
  = \norm{u-U}_{E}\,\norm{u-v}_{E},
\end{equation}
so $\norm{u-U}_{E}^{2}\le \norm{u-U}_{E}\norm{u-v}_{E}$, and cancelling one
factor:
\begin{equation}\label{eq:cea1}
  \boxed{\;\norm{u-U}_{E} \;\le\; \norm{u-v}_{E} \qquad \forall\, v\in\Vh\;}
\end{equation}
The FE solution $U$ is the best solution, meaning \emph{smallest in the energy
norm}.

\todoq{The middle line of my notes mixes the two norms
($\norm{(u-U)'}_{E}$ where it should be $\norm{u-U}_{E}$, since
$\norm{v}_{E}=\norm{v'}_{a}$ already carries the derivative). Restored above.}

\begin{intuition}[C\'ea's lemma with constant $1$]
Lecture 1's sample stated C\'ea with a factor $M/\alpha$. Here it is
\emph{exactly} $1$ --- because the energy norm is built from the same $a$ as the
bilinear form, so continuity and coercivity constants are both $1$ by
construction.

Read \eqref{eq:cea1} directly: $U$ is the closest point of $\Vh$ to $u$ in
energy. No other piecewise linear on this mesh stores a more accurate energy.
Combined with Galerkin orthogonality, $U$ is the orthogonal projection of $u$
onto $\Vh$, and Pythagoras gives $\norm{U}_{E}\le\norm{u}_{E}$ ---
\textbf{the FE model is too stiff}, as always.

Practical reading: the solver is not the limit. Accuracy is set entirely by how
well piecewise linears can represent $u$ on this mesh. Better answers come from
a better mesh, not a better solve.
\end{intuition}
